\begin{scanpage}{0012}{viii}
\begin{center}INTRODUCTION\end{center}
In recent years, there has been a noticeable increase of interest in the diffraction of electromagnetic waves by metal bodies of complex shape. Such diffraction problems with a rigorous mathematical formulation reduce to an interpretation of the wave equation or Maxwell equations with consideration of the boundary conditions on the body's surface. However, one cannot succeed in finding solutions in the case of actual bodies of a complicated configuration. This may be done only for bodies of the simplest geometric shape --- such as an infinitely long cylinder, a sphere, a disk, etc. It turns out that the resulting solutions permit one to effectively calculate the diffraction field only under the condition that the wavelength is larger than, or comparable to, the finite dimensions of the body. In the ``quasi-optical'' case, when the wavelength is a great deal less than the dimensions of the body, the rigorous solutions usually lose their practical value, and it is necessary to add to them laborious and complicated asymptotic studies. Here, the numerical methods for the solution of boundary value problems also become ineffective. Therefore, in the theory of diffraction the approximation methods which allow one to study the diffraction of sufficiently short waves by various bodies acquire great importance.

The field scattered by a given body may be calculated approximately by means of geometric optics laws (the reflection equations, see, for example [1--3]), from the principles of Huygens-Fresnel and from the equations of Kirchhoff and Cotler [3--6].

The most common method of calculation in the quasi-optic case is the principle of Huygens-Fresnel in the formulation of Kirchhoff and Cotler --- the so-called physical optics approach. The essence of this method may be summarized as follows.
\end{scanpage}
